\section{Validation and Results}

\subsection{Simulation Environment and Setup}

The simulation was done using a laptop with 3050 series 6 GB of VRAM GPU and Cuda library support. The optimizer used was Adam cite here. Table. \ref{tab:sim_parameters} contains the independent variables used in Algorithm. \ref{alg:simulation_loop}.

\begin{table}[htbp]
    \centering
    \caption{System and Simulation Parameters}
    \label{tab:sim_parameters}
    \begin{tabular}{lcc}
        \hline
        \textbf{Parameter} & \textbf{Symbol} & \textbf{Value} \\
        \hline
        Simulation Time-Step & $\Delta t$ & $1/240$ s \\
        Gravitational Acceleration & $g$ & 9.81 m/s$^2$ \\
        Simulation Duration & $t_{sim}$ & 90 s \\
        Base Frame Mass & $m_f$ & 6.8 kg \\
        Quadrotor Mass & - & 1.5 kg \\
        Payload Mass & $m_p$ & 1.0 kg \\
        Payload Volume & $V_p$ & $125 \times 10^{-6}$ m$^3$ \\
        Total System Mass & $m_{tot}$ & 7.8 kg \\
        Base System Inertia on X axis & $J_{xx}$ & 1.713 \\
        Base System Inertia on Y axis & $J_{yy}$ & 1.713 \\
        Base System Inertia on Z axis & $J_{zz}$ & 3.413 \\
        Max Thrust per Drone & $T_{max}$ & 50 N \\
        Unknown Payload Offset & $d_{x, cog}, d_{y, cog}$ & $(-0.15, 0.15)$ m \\
        PINN Window Size & $W$ & 50 steps \\
        Optimization Frequency & $N_{opt}$ & 5 steps \\
        Gradient Steps per Cycle & $S_{grad}$ & 3 steps \\
        % Lambda Data & $\lambda_{1}$ & 100 \\
        % Lambda Translational & $\lambda_{2}$ & 1 \\
        % Lambda Rotational & $\lambda_{3}$ & 1 \\
        \hline
    \end{tabular}
\end{table}

\subsection{Parameter Convergence and Estimation}

Bahas ablation studynya disini

\begin{figure}[htbp]
    \centering
    \begin{tikzpicture}
        \begin{axis}[
            width=0.98\columnwidth, 
            height=6cm,
            xlabel={Position X-axis (m)},
            ylabel={Position Y-axis (m)},
            grid=major,
            grid style={dashed, gray!30},
            legend pos=south east,
            legend style={
                font=\footnotesize, 
                cells={anchor=east}
            },
            ymin=-6,
            axis line style={thick},
            tick style={thick}
        ]
            
            \addplot[
                color=red, 
                line width=1.5pt
            ] table [x=target_x_m, y=target_y_m, col sep=comma] {data/flight_data_figure_eight_3d_ep_1.csv};
            \addlegendentry{Target Position}
            
            \addplot[
                color=blue!60!black, 
                line width=1pt, 
                dashed 
            ] table [x=actual_x_m, y=actual_y_m, col sep=comma] {data/flight_data_figure_eight_3d_ep_0.csv};
            \addlegendentry{Actual Position}
            
        \end{axis}
    \end{tikzpicture}
    \caption{The body is able to follow the desired trajectory without crashing at $1^{st}$ episode}
    \label{fig:body_position_xy_ep_0}
\end{figure}

Operating over the receding temporal window, the network rapidly mapped the unmodeled disturbance torques to spatial mass asymmetry without requiring persistent excitation maneuvers as illustrated in Fig. \ref{fig:body_position_xz_ep_1}, the body experienced an immediate lateral drift along the positive X-axis upon takeoff due to the off-center payload. However, designed controller successfully arrested this transient drift and recovered the trajectory. This rapid stabilization was enabled by the PINN periodically converging on the offset parameters (Figs. \ref{fig:offset_x_convergence_ep_1} and \ref{fig:offset_y_convergence_ep_1}) and feeding them back into the control allocation matrix. While this represents a partial parameter convergence, it proved entirely sufficient to restore system stability. By periodically identifying and compensating for the majority of the disturbance torque, the updated allocation matrix effectively restored the baseline control, allowing it to reject the remaining unmodeled dynamics without saturating the rotors.

Furthermore, the network successfully converged upon the payload mass prediction. However, as observed in Fig. (\ref{fig:mass_ep_1}), the mass estimation exhibited cyclical oscillations after the 20-second mark. These periodic fluctuations correspond directly to the dynamic tracking maneuvers, indicating that the network's mass estimation remains coupled to the multi-agent system's changing inertial movements during maneuvers.

\begin{table}[htbp]
    \centering
    \caption{Performance Metrics Across Multiple Episodes}
    \label{tab:ep_res_all}
    \begin{tabular}{lccc}
        \hline
        \textbf{Quality} & \textbf{Ep. 1} & \textbf{Ep. 5} & \textbf{Ep. 10} \\
        \hline
            RMSE Total Loss & 0.5462 & 1.5972 & 2.2172 \\
            RMSE Data Loss & 0.0023 & 0.0017 & 0.0176 \\
            RMSE Physics Loss & 0.3155 & 1.4201 & 2.042 \\
            Offset X RMSE (meters) & 0.0153 & 0.0478 & 0.0518 \\
            Offset Y RMSE (meters) & 0.0470 & 0.0762 & 0.0543 \\ 
            Trajectory RMSE (meters) & 0.4129 & 0.3757 & 0.3643 \\ 
            Trajectory x-axis RMSE (meters) & 0.2443 & 0.1521 & 0.1198 \\ 
            Trajectory y-axis RMSE (meters) & 0.2799 & 0.2946 & 0.2954 \\ 
            Trajectory z-axis RMSE (meters) & 0.1802 & 0.1768 & 0.1763 \\ 
            Avg PINN Compute Time (miliseconds) & 37.59 & 34.96 & 37.18 \\ 
            Avg PINN Compute Time (Hz) & 26.6 & 28.6 & 26.9 \\ 
        \hline
    \end{tabular}
\end{table}

\begin{table}[htbp]
    \centering
    \caption{$2^{nd}$ Episode Results}
    \label{tab:ep_res_1}
    \begin{tabular}{lcc}
        \hline
        \textbf{Quality} & \textbf{Value} & \textbf{Unit} \\
        \hline
            RMSE Total Loss & 0.6862 & - \\
            RMSE Data Loss & 0.0016 & - \\
            RMSE Physics Loss & 0.5230 & - \\
            Offset X RMSE & 0.1564 & meters \\
            Offset Y RMSE & 0.0901 & meters \\
            Trajectory Loss RMSE & 0.4129 & meters \\ 
            Trajectory x-axis RMSE & 0.2443 & meters \\ 
            Trajectory y-axis RMSE & 0.2799 & meters \\ 
            Trajectory z-axis RMSE & 0.1802 & meters \\ 
            Avg PINN Compute Time & 37.59 & ms \\ 
            Avg PINN Compute Time & 26.6 & Hz \\ 
        \hline
    \end{tabular}
\end{table}

\begin{figure}[htbp]
    \centering
    \begin{tikzpicture}
        \begin{axis}[
            width=0.98\columnwidth, 
            height=5cm,
            xlabel={Time (s)},
            ylabel={Mass (kg)},
            xtick distance=10,
            grid=major,
            grid style={dashed, gray!30},
            legend pos=south east,
            legend style={
                font=\footnotesize, 
                cells={anchor=east}
            },
            axis line style={thick},
            tick style={thick}
        ]
            
            \addplot[
                color=red, 
                line width=1.5pt
            ] table [x=time_sec, y=true_mass_kg, col sep=comma] {data/flight_data_figure_eight_3d_ep_1.csv};
            \addlegendentry{True total mass}
            
            \addplot[
                color=blue!60!black, 
                line width=1pt, 
                dashed 
            ] table [x=time_sec, y=pred_mass_kg, col sep=comma] {data/flight_data_figure_eight_3d_ep_1.csv};
            \addlegendentry{Predicted mass}
            
        \end{axis}
    \end{tikzpicture}
    \caption{Mass prediction over-time at $2^{st}$ episode. The prediction oscillates in a repeated manner following the figure 8 trajectory.}
    \label{fig:mass_ep_1}
\end{figure}

\begin{figure}[htbp]
    \centering
    \begin{tikzpicture}
        \begin{axis}[
            width=0.98\columnwidth, 
            height=5cm,
            xlabel={Time (s)},
            ylabel={Loss},
            xtick distance=10,
            grid=major,
            grid style={dashed, gray!30},
            legend pos=north east,
            legend style={
                font=\footnotesize, 
                cells={anchor=east}
            },
            axis line style={thick},
            tick style={thick}
        ]
            
            \addplot[
                color=red, 
                line width=1.5pt
            ] table [x=time_sec, y=physics_loss, col sep=comma] {data/flight_data_figure_eight_3d_ep_1.csv};
            \addlegendentry{Physics loss}
            
            % \addplot[
            %     color=blue!60!black, 
            %     line width=1pt, 
            %     dashed 
            % ] table [x=actual_x_m, y=actual_z_m, col sep=comma] {data/flight_data_figure_eight_3d_ep_1.csv};
            % \addlegendentry{Actual Position}
            
        \end{axis}
    \end{tikzpicture}
    \caption{Physics loss over-time at $2^{st}$ episode. The PINN is able to arrest the sudden loss spikes and the simulation is done without crashing.}
    \label{fig:physics_loss_ep_1}
\end{figure}

\begin{figure}[htbp]
    \centering
    \begin{tikzpicture}
        \begin{axis}[
            width=0.98\columnwidth, 
            height=5cm,
            xlabel={Position X-axis (m)},
            ylabel={Position Z-axis (m)},
            grid=major,
            grid style={dashed, gray!30},
            legend pos=south west,
            legend style={
                font=\footnotesize, 
                cells={anchor=west}
            },
            axis line style={thick},
            tick style={thick}
        ]
            
            \addplot[
                color=red, 
                line width=1.5pt
            ] table [x=target_x_m, y=target_z_m, col sep=comma] {data/flight_data_figure_eight_3d_ep_1.csv};
            \addlegendentry{Target Position}
            
            \addplot[
                color=blue!60!black, 
                line width=1pt, 
                dashed 
            ] table [x=actual_x_m, y=actual_z_m, col sep=comma] {data/flight_data_figure_eight_3d_ep_1.csv};
            \addlegendentry{Actual Position}
            
        \end{axis}
    \end{tikzpicture}
    \caption{Body frame position on XZ-axis at $2^{st}$ episode. The body immediately slides due to the off-center payload. The PINN is able to estimate the offset quickly before it reached the distance error limit.}
    \label{fig:body_position_xz_ep_1}
\end{figure}

\begin{figure}[htbp]
    \centering
    \begin{tikzpicture}
        \begin{axis}[
            width=0.98\columnwidth, 
            height=4cm,
            xlabel={Time (s)},
            ylabel={Offset (m)},
            grid=major,
            grid style={dashed, gray!30},
            legend pos=north east,
            legend style={
                font=\footnotesize, 
                cells={anchor=east}
            },
            xmin=-5,xmax=90,
            ymin=-.6,
            xtick distance=10,
            axis line style={thick},
            tick style={thick}
        ]
            
            \addplot[
                color=red, 
                line width=1.5pt
            ] table [x=time_sec, y=true_offset_x_m, col sep=comma] {data/flight_data_figure_eight_3d_ep_1.csv};
            \addlegendentry{True X offset}
            
            \addplot[
                color=blue!60!black, 
                line width=1pt, 
                dashed 
            ] table [x=time_sec, y=pred_offset_x_m, col sep=comma] {data/flight_data_figure_eight_3d_ep_1.csv};
            \addlegendentry{Predicted X offset}
            
        \end{axis}
    \end{tikzpicture}
    \caption{Predicted payload offset over time along the X-axis during the $2^{nd}$ episode. The network's ability to rapidly converge upon the true offset during the initial takeoff transient demonstrates converging parameter estimation without the requirement for persistent excitation.}
    \label{fig:offset_x_convergence_ep_1}
\end{figure}

\begin{figure}[htbp]
    \centering
    \begin{tikzpicture}
        \begin{axis}[
            width=0.98\columnwidth, 
            height=4cm,
            xlabel={Time (s)},
            ylabel={Offset (m)},
            grid=major,
            grid style={dashed, gray!30},
            legend pos=north east,
            legend style={
                font=\footnotesize, 
                cells={anchor=east}
            },
            xmin=-5,xmax=90,
            ymin=-.6,
            xtick distance=10,
            axis line style={thick},
            tick style={thick}
        ]
            
            \addplot[
                color=red, 
                line width=1.5pt
            ] table [x=time_sec, y=true_offset_y_m, col sep=comma] {data/flight_data_figure_eight_3d_ep_1.csv};
            \addlegendentry{True Y offset}
            
            \addplot[
                color=blue!60!black, 
                line width=1pt, 
                dashed 
            ] table [x=time_sec, y=pred_offset_y_m, col sep=comma] {data/flight_data_figure_eight_3d_ep_1.csv};
            \addlegendentry{Predicted Y offset}
            
        \end{axis}
    \end{tikzpicture}
    \caption{Predicted payload offset over-time on Y-axis at $2^{nd}$ episode. The network's ability to rapidly converge upon the true offset during the initial takeoff transient demonstrates converging parameter estimation without the requirement for persistent excitation.}
    \label{fig:offset_y_convergence_ep_1}
\end{figure}

\begin{figure}[htbp]
    \centering
    \begin{tikzpicture}
        \begin{axis}[
            width=0.98\columnwidth, 
            height=6cm,
            xlabel={Position X-axis (m)},
            ylabel={Position Y-axis (m)},
            grid=major,
            grid style={dashed, gray!30},
            legend pos=south east,
            legend style={
                font=\footnotesize, 
                cells={anchor=east}
            },
            ymin=-6,
            axis line style={thick},
            tick style={thick}
        ]
            
            \addplot[
                color=red, 
                line width=1.5pt
            ] table [x=target_x_m, y=target_y_m, col sep=comma] {data/flight_data_figure_eight_3d_ep_1.csv};
            \addlegendentry{Target Position}
            
            \addplot[
                color=blue!60!black, 
                line width=1pt, 
                dashed 
            ] table [x=actual_x_m, y=actual_y_m, col sep=comma] {data/flight_data_figure_eight_3d_ep_1.csv};
            \addlegendentry{Actual Position}
            
        \end{axis}
    \end{tikzpicture}
    \caption{Target position and actual position at $2^{nd}$ episode. Even with a sudden drift during takeoff, the controller is able to stabilize and follow the figure 8 trajectory.}
    \label{fig:body_position_xy_ep_1}
\end{figure}